% Ad Familiares 13.33 --- headnote
% ad-familiares-13-33, 46 BC, Rome

\textit{Cicero to Manius Acilius Glabrio, proconsul of
Sicily, written from Rome in 46 BC (the manuscript
dateline: \textit{Scr. Romae, ut videtur, a. 708 (46)}).
The second of the surviving cluster of recommendations
to Acilius (Fam.\ 13.30--39). The beneficiary is Gnaeus
Otacilius Naso --- a Roman businessman of equestrian rank,
otherwise obscure, with commercial interests in Sicily
managed on the ground by his freedmen Hilarus, Antigonus,
and Demostratus. Cicero stands in close personal regard
of him at Rome and asks the proconsul to extend the same
favour to Naso's affairs and to his agents.}

\textit{The letter is among the most compressed of the
\textit{commendaticiae}: a single paragraph that gestures
briefly at the sketch of the man (humanity, uprightness,
daily intimacy), waives the formal rehearsal of his
merits as already implied by that intimacy, names the
agents who stand in for him in the province, and closes
with the standard request that the recommendation be
felt to have weighed with the addressee. The economy is
itself the point. The genre had become a set form, and
Cicero handles its formulae with a craftsman's lightness,
varying the touch from letter to letter without ever
straying from the shape.}
